\FloatBarrier
\subsection{Exponential Runtime Growth Analysis of the Exact Algorithm}
\label{sec:exact_dp_runtime_analysis}

During the execution of the exact dynamic programming algorithm from Section~\ref{sec:log_n_log_log_n_approx}, it could be noticed that the running time of the procedure behaved better then $O\br{2^nn}$. We decided to find the value of $b$ such that the function $T(n) = b^n$ best approximates the running time of the procedure. 
The coefficient of determination $R^2$ indicates the quality of the exponential fit. It is defined as:
$$
R^2 = 1 - \frac{\sum_{i=1}^{n}\br{y_i - \hat{y}_i}^2}{\sum_{i=1}^{n}\br{y_i - \bar{y}}^2}
$$
where $y_i$ are the observed runtime values, $\hat{y}_i$ are the predicted values from the model $T(n) = b^n$, and $\bar{y}$ is the mean of observed values. Values of $R^2$ close to 1 suggest that the exponential model accurately describes the runtime growth, while values close to 0 indicate poor fit.  We calculated the values of $b$ for which the value of $R^2$ is maximized:
\FloatBarrier

\begin{table}[H]
\centering
\caption{Exponential fit for exact DP algorithm runtime}
\label{tab:dp_exponential_fits}
\small
\begin{tabular}{|l|l|c|c|c|c|}
\hline
\multirow{2}{*}{\textbf{Variant}} & \multirow{2}{*}{\textbf{Distribution}} & \multicolumn{2}{c|}{\textbf{Average CPU Time}} & \multicolumn{2}{c|}{\textbf{Maximum CPU Time}} \\
\cline{3-6}
 &  & $b$ & $R^2$ & $b$ & $R^2$ \\
\hline
\multirow{3}{*}{Random} & Binomial & 1.341 & 0.924 & 1.381 & 0.734 \\
 & Exponential & 1.348 & 0.606 & 1.382 & 0.465 \\
 & Uniform & 1.350 & 0.918 & 1.386 & 0.598 \\
\hline
\multirow{3}{*}{Random Degree 3} & Binomial & 1.355 & 0.918 & 1.377 & 0.961 \\
 & Exponential & 1.358 & 0.978 & 1.379 & 0.940 \\
 & Uniform & 1.353 & 0.975 & 1.373 & 0.967 \\
\hline
\multirow{3}{*}{Random Diameter 6} & Binomial & 1.518 & 0.817 & 1.600 & 0.251 \\
 & Exponential & 1.519 & 0.717 & 1.600 & 0.764 \\
 & Uniform & 1.509 & 0.882 & 1.610 & 0.949 \\
\hline
\end{tabular}
\end{table}
\FloatBarrier
We observe that:
\begin{itemize}
    \item The bases $b$ are all significantly smaller than 2, ranging from approximately 1.34 to 1.61, indicating that the algorithm performs better than the worst-case $O(2^n)$ growth in practice for these instances.
    \item Random trees with bounded diameter (Diameter 6) display the highest growth bases ($b \approx 1.51$-$1.61$).
    \item Other tree classes exhibited value of $b$ oscillating around $b \approx 1.34$-$1.39$ which is siginificantly better then the worst case $O\br{2^nn}$ running time.
    \item For each class of trees the values of $b$ are relatively consistent across different cost distributions (Uniform, Binomial, Exponential), suggesting that the distribution of query costs has a limited impact on the exponential growth rate of the algorithm's runtime and that perhaps the values of $b$ might stabilize for larger input sizes.
\end{itemize}

\FloatBarrier
